For a graph $G$:
\begin{itemize}
\item The \emph{chromatic number} $\chi(G)$ is the minimum $m$ such that $V(G)$ can be partitioned into $m$ independent sets.
\item The \emph{cochromatic number} $\zeta(G)$ is the minimum $m$ such that $V(G)$ can be partitioned into $m$ parts,
where each part induces either a clique or an independent set.
\item For an \emph{orientation} $D$ of $G$ (replace each undirected edge by one of the two possible directions),
let $\chi_d(D)$ be the minimum number of colours needed to colour $V(D)$ so that each colour class induces an
\emph{acyclic} digraph (equivalently: contains no directed cycle). This is the standard \emph{dichromatic number} of $D$.
\item Define the (undirected) \emph{dichromatic number} of $G$ as
\[
\delta(G):=\max\{\chi_d(D): D\ \text{is an orientation of }G\}.
\]
\end{itemize}

The questions are:
\begin{enumerate}
\item[(Q1)] Does $\chi(G)$ unbounded force $\delta(G)$ unbounded? Equivalently: for every $k$ does there exist $f(k)$
so that $\chi(G)\ge f(k)$ implies $\delta(G)\ge k$?
\item[(Q2)] Does $\zeta(G)$ large force the existence of a (not necessarily induced) subgraph $H\subseteq G$ with large $\delta(H)$?
\end{enumerate}
